% --------------------------------------------------------------------------------------------------------------------%
% Defining the chemical system in each geochemical modeling problem
% --------------------------------------------------------------------------------------------------------------------%
%
\section{Defining the chemical system in each geochemical modeling problem}
\begin{frame}[<+->]{Defining the chemical system for the modeling problem of interest}
\begin{itemize}
\item A chemical system is a collection of phases (one or more). Examples:
\begin{itemize}
\item a chemical system with only an aqueous phase;
\item a chemical system with one aqueous phase and one gaseous phase;
\item a chemical system with aqueous, gaseous, and mineral phases.
\end{itemize}
\item Each phase has one or more chemical species:
\begin{itemize}
\item aqueous: H$_{2}$O(aq), H$^{+}$(aq), OH$^{-}$(aq), HCO$_{3}^{-}$(aq),
CO$_{3}^{2-}$(aq), CO$_{2}$(aq), Ca$^{2+}$(aq)
\item gaseous: H$_{2}$O(g), CO$_{2}$(g), O$_{2}$(g), H$_{2}$S(g)
\item calcite: CaCO$_{3}$(s)
\item halite: NaCl(s)
\end{itemize}
\item In geochemical modeling, suitable phases and their species need to
be considered. Either manually, or automatically from databases.
\end{itemize}
\end{frame}
%
% --------------------------------------------------------------------------------------------------------------------%
% Example of chemical system when substances H$_{2}$O and CO$_{2}$
% are involved
% --------------------------------------------------------------------------------------------------------------------%
%
\begin{frame}{Example of chemical system when substances H$_{2}$O and CO$_{2}$
are involved}

\lcol
\begin{itemize}
\item Assume water (H$_{2}$O) and carbon dioxide (CO$_{2}$), are mixed
at some given temperature $T$ and pressure $P$. 
\item To model the equilibrium state resulting from this process, the first
step is to think about:
\begin{itemize}
\item the possible phases of matter that could exist; and
\item the species in each phase. 
\end{itemize}
\end{itemize}
\rcol
\begin{itemize}
\item A simple (but reasonable) chemical system for this problem is:
\end{itemize}
\begin{center}
\begin{tabular}{cc}
\toprule 
\textbf{Aqueous Phase} & \textbf{Gaseous Phase}\tabularnewline
\midrule
H$_{2}$O(aq) & H$_{2}$O(g)\tabularnewline
H$^{+}$(aq) & CO$_{2}$(g)\tabularnewline
OH$^{-}$(aq) & \tabularnewline
HCO$_{3}^{-}$(aq) & \tabularnewline
CO$_{3}^{2-}$(aq) & \tabularnewline
CO$_{2}$(aq) & \tabularnewline
\bottomrule
\end{tabular}
\par\end{center}

\ecol
\end{frame}
%
% --------------------------------------------------------------------------------------------------------------------%
% Example of chemical system when substances H2O and CaCO3 are involved
% --------------------------------------------------------------------------------------------------------------------%
%
\begin{frame}{Example of chemical system when substances H$_{2}$O and CaCO$_{\boldsymbol{3}}$
are involved}
\begin{itemize}
\item Assume water (H$_{2}$O) and calcite (CaCO$_{3}$) are mixed at some
given temperature $T$ and pressure $P$. 
\item A simple (but reasonable) chemical system for this problem is:
\end{itemize}
\begin{center}
\begin{tabular}{cc}
\toprule 
\multicolumn{1}{c}{\textbf{Aqueous Phase}} & \textbf{Calcite Phase}\tabularnewline
\midrule
H$_{2}$O(aq) & CaCO$_{3}$(s, calcite)\tabularnewline
H$^{+}$(aq) & \tabularnewline
OH$^{-}$(aq) & \tabularnewline
HCO$_{3}^{-}$(aq) & \tabularnewline
CO$_{3}^{2-}$(aq) & \tabularnewline
CO$_{2}$(aq) & \tabularnewline
Ca$^{2+}$(aq)  & \tabularnewline
\bottomrule
\end{tabular}
\par\end{center}

\end{frame}
%
% --------------------------------------------------------------------------------------------------------------------%
% Example of chemical system when substances H2O, NaCl and CO2 are involved
% --------------------------------------------------------------------------------------------------------------------%
%
\begin{frame}{Example of chemical system when substances H$_{2}$O, NaCl and CO$_{2}$
are involved}
\begin{itemize}
\item Assume water (H$_{2}$O), carbon dioxide (CO$_{2}$), and sodium chloride
(NaCl), are mixed at some given temperature $T$ and pressure $P$. 
\item A simple (but reasonable) chemical system for this problem is:
\end{itemize}
\begin{center}
\begin{tabular}{cccc}
\toprule 
\multicolumn{2}{c}{\textbf{Aqueous Phase}} & \multicolumn{1}{c}{\textbf{Gaseous Phase}} & \textbf{Halite Phase}\tabularnewline
\midrule
H$_{2}$O(aq) & CO$_{2}$(aq) & CO$_{2}$(g) & NaCl(s, halite)\tabularnewline
H$^{+}$(aq) & HCO$_{3}^{-}$(aq) & H$_{2}$O(g) & \tabularnewline
OH$^{-}$(aq) & CO$_{3}^{2-}$(aq) &  & \tabularnewline
Na$^{+}$(aq) &  &  & \tabularnewline
Cl$^{-}$(aq) &  &  & \tabularnewline
NaCl(aq) &  &  & \tabularnewline
\bottomrule
\end{tabular}
\par\end{center}

\end{frame}
%
% --------------------------------------------------------------------------------------------------------------------%
% How many phases and species to consider for our chemical system?
% --------------------------------------------------------------------------------------------------------------------%
%
\begin{frame}{How many phases and species to consider for our chemical system?}

\lcol

\small
\begin{itemize}
\item In general, the \textbf{number of phases and species} should be \textbf{as
large as possible} when \emph{defining a chemical system}.
\item However, the calculations get more expensive with increasing number
of species and phases.
\item \textbf{\alert{\textbf{Important:}}} Not all phases considered in
the calculation will actually exist in positive amounts. Their existence
depends on the input conditions.
\item The ternary phase diagram on the right shows conditions in which not
all phases are \textbf{stable at equilibrium}.
\end{itemize}
\rcol
\begin{center}
\includegraphics[width=0.9\columnwidth]{figures/ternary-phase-diagram-h2o-co2-nacl}
\par\end{center}

\ecol
\end{frame}
%
% --------------------------------------------------------------------------------------------------------------------%
% Final considerations about defining the chemical system
% --------------------------------------------------------------------------------------------------------------------%
%
\begin{frame}[<+->]{Final considerations about defining the chemical system}
\begin{itemize}
\item In most computer codes for modeling geochemical reactions, no manual
selection of phases and species are needed. 
\item These can be determined automatically by searching in \textbf{thermodynamic
databases} all possible species and phases that could exist for a
given model input. 
\end{itemize}
\end{frame}
